% 05_discussion.tex -- discussion and conclusion, draft 23 Oct 2026.
\section{Discussion}
\label{sec:discussion}

\paragraph{Rigid links in the twin, slack tethers on the rig}
The twin models every cable as a rigid link, which can push as well as pull and welds wherever its end lies. This
flatters the rejoin and also creates its own failure: on a T, the rejoiner's rigid link set the ring's pitch, so a
rising height reference lifted the ring's free side, and the height hold that cured it is a patch for the twin's
links. On the rig the tethers go slack, so that lift cannot arise, but the attach is harder: the magnet catches only on
contact, which is why the rig's attach is confirmed from motion capture before the newcomer is folded in. The rig's
release reaction was measured on the rig itself; the twin's release on 1\,m cables gave a similar tilt peak but a
slower return to band, and the rig flew 0.55\,m rods, so the two are not a like-for-like validation.

\paragraph{Timing}
Every reconfiguration in Section~\ref{sec:results} is paced by detection and blending rather than by the solver: a
lost carrier is seen from its cable's stretch in 0.33--0.42\,s and the problem is resized within about 0.1\,s more,
and a rejoin's fold-in is blended over 10\,s. The solver sets the size of the team instead. At four carriers one SQP
iteration per tick flies as well as five in the twin's closed loop, but its plans are not converged on the rig's
flights, and at eight carriers only five iterations pass every bar (Section~\ref{sec:res_scale}). The eight-carrier
chain ran slower than real time, so it shows that the plan and the physics compose, not that the rig's laptop can fly
it.

\paragraph{Geometry sets how hard a release hits}
The release transient is set less by the controller than by the team that remains. The ring's tilt after a release
grows steeply with the largest gap the survivors leave around the rim, and on 1\,m cables a 165$^\circ$ gap already
crosses the 30$^\circ$ drop line (Section~\ref{sec:res_release_gap}). The planner can shorten the transient, most of
all when it plans from the measured cable pull rather than the planned one, but it cannot remove the torque that an
uneven ring of survivors applies. A team that must shed carriers should therefore leave in an order that keeps the
survivors' gaps small, and a drop rule written in terms of that gap, rather than of the number of carriers left,
states what the ring can bear. On the attach side, folding in at the measured catch point makes the twin's fold-in
tolerant of catch-point error, 16\,cm radially and 6\,cm along the rim (Section~\ref{sec:res_attach_offset}); past
about 2.5\,cm this is tolerance to model error, not a catch a real magnet could make.

\paragraph{Limitations}
The rig results cover the release only, three runs on 0.55\,m rods; the rejoin, the attach-offset sweep, the gap
sweep and the eight-carrier chain are twin results, two runs a cell or fewer, on rigid links that weld wherever the
end lies. No reconfiguration is compared with a baseline that does not resize the problem. The eight-carrier chain ran
slower than real time at one iteration a tick, below the five that the offline check asks for at that size, and
along the rim the measured catch point's advantage over the slot assumed is not shown.

\paragraph{Next on the rig}
The next rig flights test the attach of a newcomer on 1\,m cables, three carriers to four with a basket, with the
attach confirmed by motion capture before the fold-in; the four-carrier hover and paths on 1\,m cables; and the
planner's tick in flight with the speed-ups, which sets the real-time limit. The rejoin of a departed carrier, which
in the twin depends on the joint's weld, waits for a rig attach that has been shown repeatable.

\section{Conclusion}
\label{sec:conclusion}

We presented a planner that lets a team of quadrotors carry a cable-suspended ring through changes of its own
membership in flight. On the rig, an unannounced release of one of four carriers was detected in 0.42\,s, the team
was resized 0.09\,s later, and the ring tilted at most 6.5$^\circ$ before returning to its band in 1.8\,s. In the twin,
a departed carrier flew back to its own plate, welded within 1.5\,cm and was folded in with the ring tilting
1.3--3.4$^\circ$ in 10 of 11 runs, and in one exploratory run eight carriers took a larger ring down to four and back with a tilt peak of 3.35$^\circ$. In offline
replays of nine flights, two SQP iterations per tick converge for every team of up to four, and the iterations needed grow with the size of the team
(Section~\ref{sec:res_scale}). How hard a release hits is set by the survivors' geometry, the largest gap they leave,
and folding a newcomer in at its measured catch point tolerates large errors in that point. The rig attach of a newcomer on 1\,m cables is the
next step \placeholder{result from 29 Oct}, together with the real-time limit of the planner's tick on the rig
\placeholder{real-time limit}.
